%/*% pdflatex -file-line-error -interaction nonstopmode %

\documentclass{article}
\usepackage{graphicx}
\usepackage{amsmath}
\usepackage{amsfonts}
\usepackage{hyperref}
\usepackage{exercises}
\usepackage{amsthm}
\usepackage{yhmath}
\newcommand*{\oldneg}{\mathord{\sim}}

\theoremstyle{definition}
\newtheorem{definition}{Definition}

\setlength{\parindent}{.5in}
\setlength{\parskip}{12pt}

\title{1.5 Exercises}
\begin{document}
\maketitle

\section*{1}
The radicand.

\section*{2}
Consider some fraction with a denominator that is a square root, or is of the form $b+\sqrt{a},\ b-\sqrt{a}$. To fully simplify this fraction, one needs 
to remove any radicals from that denominator, which is done through a proccess called rationalizing the denominator. In the first case, this can be done 
by multiplying the fraction by $\frac{\sqrt{b}}{\sqrt{b}}$, and it can be done in the second case by multiplying the fraction by $\frac{a-\sqrt{b}}{a-\sqrt{
b}}$, and it can be done in the third case by multiplying the fraction by $\frac{a+\sqrt{b}}{a+\sqrt{b}}$.

\section*{3}
$\sqrt{28}=\sqrt{7}\times \sqrt{4}\newline =2\sqrt{7}$

\section*{4}
$\sqrt{192}=\sqrt{64}\times \sqrt{3}\newline =8\sqrt{3}$

\section*{5}
$\sqrt{150}=\sqrt{25}\times \sqrt{6}\newline =5\sqrt{6}$

\section*{6}
$\sqrt{3}\times \sqrt{27}=\sqrt{81}=9$

\section*{7}
$4\sqrt{6}\times \sqrt{6}=4\times 6=24$

\section*{8}
$5\sqrt{24}\times 3\sqrt{10}=15\sqrt{8}\sqrt{3}\sqrt{5}\sqrt{2}=15\sqrt{16}\times\sqrt {15}=60\sqrt{15}$

\section*{9}
$\sqrt\frac{5}{16}=\frac{\sqrt{5}}{\sqrt{16}}=\frac{\sqrt{5}}{4}$

\section*{10}
$\sqrt\frac{35}{36}=\frac{\sqrt{35}}{\sqrt{36}}=\frac{\sqrt{35}}{6}$

\section*{11}
$\frac{8}{\sqrt{3}}=\frac{8\times \sqrt{3}}{\sqrt{3}\sqrt{3}}=\frac{8\sqrt{3}}{3}$

\section*{12}
$\frac{7}{\sqrt{12}}=\frac{7\sqrt{12}}{\sqrt{12}\sqrt{12}}=\frac{7\sqrt{12}}{12}$

\section*{13}
$\sqrt\frac{18}{11}=\frac{\sqrt{18}}{\sqrt{11}}=\frac{3\sqrt{2}\sqrt{11}}{\sqrt{11}\sqrt{11}}=\frac{3\sqrt{22}}{11}$

\section*{14}
$\sqrt\frac{13}{28}=\frac{\sqrt{13}}{\sqrt{28}}=\frac{\sqrt{13}}{2\sqrt{7}}=\frac{\sqrt{13}\sqrt{7}}{2\sqrt{7}\sqrt{7}}=
\frac{\sqrt{91}}{14}$

\section*{15}
$\frac{2}{1-\sqrt{3}}=\frac{2\times (1+\sqrt{3})}{(1-\sqrt{3})\times(1+\sqrt{3})}\newline 
=\frac{2+2\sqrt{3}}{-2}=\frac{1+\sqrt{3}}{-1}$

\section*{16}
$\frac{1}{5+\sqrt{6}}=\frac{1\times(5-\sqrt{6})}{(5+\sqrt{6})\times(5-\sqrt{6})}=\frac{5-\sqrt{6}}{25-6}=
\frac{5-\sqrt{6}}{19}$

\section*{17}
$\frac{\sqrt{2}}{4+\sqrt{5}}=\frac{\sqrt{2}\times(4-\sqrt{5})}{(4+\sqrt{5})\times (4-\sqrt{5})}=
\frac{4\sqrt{2}-\sqrt{10}}{11}$

\section*{18}
$\frac{3+\sqrt{7}}{2-\sqrt{10}}=\frac{3+\sqrt{7}\times(2+\sqrt{10})}{(2-\sqrt{10})\times (2+\sqrt{10})}=
\frac{6+2\sqrt{7}+3\sqrt{10}+\sqrt{70}}{-6}$

\section*{19}
$\sqrt{108}=\sqrt{36}\sqrt{3}=6\sqrt{3}$, and thus, the correct answer is $C$.

\section*{20}
The answer does not factor out the perfect square factors of $24$:\newline
$\sqrt{96}=\sqrt{16}\sqrt{6}=4\sqrt{6}$

\section*{21}
The error is that the ngeative solution to the quadratic equation was ignored. Instead, the answer is that $x=9\vee x=-9$.

\section*{22}
$s^2=169\newline\rightarrow s=\pm 13$

\section*{23}
$a^2=50\newline\rightarrow a=\pm\sqrt{50}=\pm 5\sqrt{2}$

\section*{24}
$x^2=84\newline\rightarrow x=\pm\sqrt{84}=\pm 2\sqrt{21}$

\section*{25}
$6z^2=150\newline\rightarrow z^2=25\newline\rightarrow z=\pm 5$

\section*{26}
$4p^2=448\newline\rightarrow p^2=112\newline\rightarrow p=\pm \sqrt{112}=\pm 4\sqrt{7}$

\section*{27}
$-3w^2=-213\newline\rightarrow w^2=71\newline\rightarrow w=\pm \sqrt{71}$

\section*{28}
$7r^2-10=25\newline\rightarrow r^2=5\newline\rightarrow r=\pm \sqrt{5}$

\section*{29}
$\frac{x^2}{25}-6=-2\newline\rightarrow x^2=100\newline\rightarrow x=\pm 10$

\section*{30}
$\frac{t^2}{20}+8=15\newline\rightarrow t^2=140\newline\rightarrow t=\pm\sqrt{140}\newline\rightarrow t=\pm 2\sqrt{35}$

\section*{31}
$4(x-1)^2=8\newline\rightarrow (x-1)^2=2\newline\rightarrow x-1=\pm\sqrt{2}\newline\rightarrow x=1\pm\sqrt{2}$

\section*{32}
$7(x-4)^2-18=10\newline\rightarrow (x-4)^2=4\newline\rightarrow x-4=\pm 2\newline\rightarrow x=6\vee x=2$

\section*{33}
$2(x+2)^2-5=8\newline\rightarrow (x+2)^2=\frac{13}{2}\newline\rightarrow x+2=\pm\sqrt\frac{13}{2}\newline\rightarrow x=\sqrt\frac{13}{2}-2\vee 
x=-\sqrt\frac{13}{2}-2$

\section*{34}
(a): $\sqrt{12}-x\sqrt{2}\le \sqrt{3}\newline\rightarrow x\sqrt{2}\le \sqrt{3}\newline\rightarrow x\ge \frac{\sqrt{3}}{\sqrt{2}}\newline\rightarrow 
x\ge \sqrt\frac{3}{2}$\newline
(b): $x\sqrt{3}>2+x\newline\rightarrow x\sqrt{3}-x>2\newline\rightarrow (\sqrt{3}-1)x>2\newline\rightarrow x>\frac{2}{-1+\sqrt{3}}\newline\rightarrow 
x>\frac{2(-1-\sqrt{3})}{-2}\newline\rightarrow x>-(-1-\sqrt{3})\newline\rightarrow x>1+\sqrt{3}$\newline
(c): $2\sqrt{5}=4+x\sqrt{8}\newline\rightarrow 4+2\sqrt{2}x=2\sqrt{5}\newline\rightarrow 2\sqrt{2}x=2\sqrt{5}-4\newline\rightarrow 
x=\frac{2\sqrt{5}-4}{2\sqrt{2}}\newline\rightarrow =\frac{4\sqrt{10}-8\sqrt{2}}{8}\newline\rightarrow =\frac{\sqrt{10}-2\sqrt{2}}{2}$

\section*{35}
One method is two move $4$ to the right side of the equation, and then take its square roots, thus yielding $x=2\vee x=-2$. Alternatively, we 
can factor this equation by using the differences of squares pattern:\newline $(x+2)(x-2)=0\newline\rightarrow x=2\vee x=-2$

\section*{36}
(a): $x^2=9\newline\rightarrow x=\pm 3$\newline
(b): $x^2=0\newline\rightarrow x=0$\newline
(c): $x^2=-4\newline\rightarrow x=\sqrt{-4}$

\section*{37}
$a(x+b)^2=c\newline\rightarrow (x+b)^2=\frac{c}{a}\newline\rightarrow x+b=\pm \sqrt\frac{c}{a}\newline\rightarrow 
x+b=\pm \frac{\sqrt{c}}{\sqrt{a}}\newline\rightarrow x+b=\pm \frac{\sqrt{ac}}{a}\newline\rightarrow x=\pm \frac{\sqrt{ac}}{a}-b$

\section*{38}
$h=-16t^2+40\newline\rightarrow 0=-16t^2+40\newline\rightarrow t^2=\frac{5}{2}\newline\rightarrow t=\pm \sqrt\frac{5}{2}\approx\pm 1.58$
, and thus, the diver is in 
the air for about $1.58$ seconds. 

\section*{39}
Earth: $0=-\frac{32}{2}t^2+150\newline\rightarrow t^2=9.375\newline\rightarrow t=\pm \sqrt{9.375}\approx\pm 3.062$, meaning it takes about $3.062$ seconds 
for the rock to fall on Earth.\newline
Mars: $0=-\frac{12}{2}t^2+150\newline\rightarrow t^2=25\newline\rightarrow t=\pm 5$, meaning it takes $5$ seconds for the rock to fall on Mars.\newline 
Jupiter: $0=-\frac{76}{2}t^2+150\newline\rightarrow t^2=\frac{75}{19}\newline\rightarrow t=\pm\sqrt\frac{75}{19}\approx \pm 1.97$, and thus, it takes 
about $1.97$ seconds for the rock to fall on Jupiter.\newline
Saturn: $0=-\frac{30}{2}t^2+150\newline\rightarrow t^2=10\newline\rightarrow t=\pm\sqrt{10}\approx\pm 3.162$, and thus, it takes about $3.162$ seconds 
for the rock to fall on Saturn.\newline
Pluto: $0=-\frac{2}{2}t^2+150\newline\rightarrow t^2=150\newline\rightarrow t=\pm\sqrt{150}\approx\pm 12.247$, and thus, it takes about $12.247$ seconds 
for the rock to fall on Pluto.

\section*{40}
Wind for $5$ ft. wave: $5=0.019s^2\newline\rightarrow s^2=\frac{5}{0.019}\newline\rightarrow s=\pm\sqrt\frac{5}{0.019}\approx\pm 16.222$\newline
Wind for $20$ ft. wave: $20=0.019s^2\newline\rightarrow s^2=\frac{20}{0.019}\newline\rightarrow s=\pm\sqrt\frac{20}{0.019}\approx\pm 32.444$\newline
Hence, to produce a wave $20$ ft. high, it requires double the wind speed ($32.444$) as it does to produce a $5$ ft. wave ($16.222$).

\section*{41}
(a): $100 = \pi r^2$\newline
(b): $r^2=\frac{100}{\pi}\newline\rightarrow r=\pm\sqrt\frac{100}{\pi}\newline\rightarrow r=\pm\frac{\sqrt{100}}{\sqrt{\pi}}=\pm\frac{10}{\sqrt{\pi}}$, 
and thus, because only positive values are accepted, $r=\frac{10}{\sqrt{\pi}}\approx 5.64$
(c): We consider any square with side lengths $s$. A circle with the same area as this square implies $s^2=\pi r^2\newline\rightarrow 
r^2=\frac{s^2}{\pi}\newline\rightarrow r=\sqrt\frac{s^2}{\pi}=\frac{s}{\sqrt{\pi}}=\frac{s\sqrt{\pi}}{\pi}$

\section*{42}
(a): $5=0.00829s^2\newline\rightarrow s^2=\frac{5}{0.00829}\newline\rightarrow s=\pm\sqrt\frac{5}{0.00829}\approx \pm 24.56$, and thus, the bicyclist 
was going about $24.56$ mph.\newline
(b): Consider a bicyclist having an air resistance of $r$ and a speed of $s$. Doubling the speed yields the folliwng equation:\newline 
$R=0.00829(2s)^2=4(0.00829s^2)$, and thus, because $r=0.00829s^2$, doubling the speed quadruples the air resistance. 

\section*{43}
To solve this, we set the given equation to $0$, and then solve for $t$:\newline
$0=(\sqrt{h_0}-\frac{2\pi d^2\sqrt{3}}{lw}t)^2\newline\rightarrow 0=\sqrt{h_0}-\frac{2\pi d^2\sqrt{3}}{lw}t\newline\rightarrow
\frac{2\pi d^2\sqrt{3}}{lw}t=\sqrt{h_0}\newline\rightarrow t=\frac{\sqrt{h_0}}{\frac{2\pi d^2\sqrt{3}}{lw}}\newline\rightarrow =
\frac{\sqrt{h_0}lw}{2\pi d^2\sqrt{3}}\newline\rightarrow =\frac{\sqrt{3h_0}lw}{6\pi d^2}$
\end{document}








